\scanpage{83}\mspage{82}

\noindent\llap{\smash{\rule[-4.9\baselineskip]{.35pt}{5.6\baselineskip}\hspace{.20em}\rule[-4.9\baselineskip]{.35pt}{5.6\baselineskip}}\hspace{.6em}}\uline{C 8}\textsubscript{(0)} Considérons les situations
\margin{\uline{provisoire !}}
\[
\begin{aligned}
\text{(a)}\qquad &X\leq\Phi\preceq X & \Phi&\in\mathcal F,\ X\in M\\
\text{(b)}\qquad &X\preceq\Psi\leq X & \deleted{\unsure{F|\Phi|\mathcal F}}\Psi&\in\mathcal F,\ X\in M
\end{aligned}
\]
\deleted{\unsure{C|A|On}}Alors on a $X=\Phi\quad\text{et}\quad X=\Psi$

Achevons la dém. Appliquant (a) à
\[
X\leq F_X\preceq X
\]
on trouve $X=F_X$, d'où, comme
\[
X\leq\bar G_X\leq F_X=X\quad,
\]
$X=\bar G_X$, d'où
\[
\deleted{\unsure{X=|X\leq|G_X=|G_X}}\bar G_X=X\preceq G_X\leq X
\]
et appliquant (b), on trouve $\deleted{\unsure{G|X|\Phi}}X=G_X$ et par suite $X\leq G_X$, qed.

\uline{Remarque} Bien sûr, C8 est nécessaire si on veut que $\ll$ soit une relation d'ordre. Et c'est aussi suffisant\unsure{,:|;:|;}

\uline{Cor} (de C8) \inserted{via prop. p\deleted{\unsure{g|.}}~76)} Si $F\ll G$ et $G\ll F$, on a $F=G$

\uline{Dém} \deleted{Soit $X\in\widetilde F$ \unsure{i.e.|ie.} $X\in M$, $X\leq F$} On va prouver $\unsure{F|\bar F}\leq G$\ ; \deleted{\unsure{(on|Car}} on aura par symétrie $G\leq F$, donc $F=G$.
\deleted{Il faut prouver \deleted{\unsure{pour|que}} pour $X\in M$, $(X\unsure{\leq|<}F)\ \unsure{-|\Rightarrow|=} X\leq G$.}

\diagram{\begin{tikzcd}[ampersand replacement=\&,column sep=0.45em,row sep=0.1em]
\deleted{F} \& \deleted{\ll} \& \deleted{G} \& \deleted{\ll} \& \deleted{F}\\
\rotatebox{90}{$\deleted{\leq}$} \& \& \rotatebox{90}{$\deleted{\leq}$} \& \& \rotatebox{90}{$\deleted{\leq}$}\\
\deleted{F_X} \& \deleted{\ll} \& \deleted{G_X} \& \deleted{\ll} \& \deleted{X}\\
\rotatebox{90}{$\deleted{\leq}$}\\
\deleted{X}
\end{tikzcd}}
